\chapter{Arrays and records}\label{ch-data}

The values of the previous chapters were single numbers, characters and
truth values. Real programs work with collections of values: a list of
measurements, a matrix, a table of names and dates. Algol 68 provides two
ways of building such values. An \term{array} holds a number of values of
the same type, selected by position; a \term{record} holds a fixed set of
values of possibly different types, selected by name. Both can be
combined freely: arrays of records, records containing arrays, arrays of
arrays. This chapter also introduces type declarations, which give a
name to a type.

\section{Arrays}\label{sec-arrays}

An array variable is declared by giving its bounds in square brackets
before the element type:

\example{ch04/array-basics}

The declaration \code|[1:5] int squares| allocates five cells for
\code|int| values, numbered 1 to 5. Unlike in C, Java or Python, the
numbering need not start at 0 or 1: \code|[0:2] char letters| has the
indices 0, 1 and 2, and an array of the years 1990 to 2025 could be
declared as \code|[1990:2025] real|. The lower bound is the first number
and the upper bound the second; the operators \code|LWB| and \code|UPB|
give them. A loop over all elements should use them rather than repeat
the numbers. If the lower bound is 1, it can be left out: \code|[n] real|
means \code|[1:n] real|.

An element is selected by an \term{index} in square brackets, as in
\code|squares[3]|. The Report calls an array a \rr{multiple value} or
\rr{row}, and the selection of an element a \rr{slice}.

The bounds are evaluated when the declaration is carried out, so the size
of an array can depend on input. The following program reads a count and
then that many numbers:

\example{ch04/dynamic-size}

\noindent The input was:
\lstinputlisting[style=output]{examples/ch04/dynamic-size.in}

Once declared, the array keeps its bounds; an array that can grow or
shrink is the subject of Chapter~\ref{ch-advanced}. Every index is checked: an index
outside the bounds stops the program with an error.

\example{ch04/out-of-bounds}

\section{Array literals}\label{sec-displays}

A list of values in parentheses, separated by commas, is an \term{array
literal} (\rr{row display}) wherever an array is expected. A constant
array can take its bounds from the literal: the declaration
\code|[] int primes = (2, 3, 5, 7, 11)| gives no bounds, and
\code|primes| has the bounds 1 to 5. The Report calls a type written
without bounds, such as \code|[] int|, a \rr{formal declarer}; one with
bounds, as in the declaration of a variable, an \rr{actual declarer}.

\example{ch04/displays}

The elements of a literal are converted to the element type where
necessary, so the integers of \code|weights| are widened to \code|real|.
The empty literal \code|()| gives an array with the bounds 1 to 0, that
is, with no elements. A literal can also be assigned to an array
variable, which replaces all its elements. A literal always has the lower
bound 1, so its bounds match those of the variable only if the variable
has the lower bound 1 and as many elements as the literal. For a variable
\code|[0:2] int w|, the assignment \code|w := (1, 2, 3)| fails, but
\code|w[@1] := (1, 2, 3)| works: \code|w[@1]| is the whole of \code|w|,
renumbered to start at 1 (Section~\ref{sec-slices}).

\section{Arrays are values}\label{sec-array-values}

An array in Algol 68 is a value, like a number. Assigning one array
variable to another copies the elements:

\example{ch04/copy}

After the assignment, \code|a| and \code|b| are separate arrays; changing
\code|b| does not change \code|a|. This is the behavior of arrays in C
when they are inside a \texttt{struct}, or of the arrays of Fortran, but
not that of Java, Python or JavaScript, where an assignment makes two
names refer to the same array. In Algol 68, two names refer to the same
array only if a reference to it is shared explicitly (Chapter~\ref{ch-advanced}).
A constant declared with an array value, \code|[] int c = b|, is a copy
as well: later changes of \code|b| do not show in \code|c|. The same holds
for a parameter of type \code|[] int| (Chapter~\ref{ch-procs}).

\begin{implnote}
Algol 68 Genie does not always copy: if a procedure has a \code|[] int|
parameter and a \code|ref [] int| parameter, and both are given the same
array, a change through the reference shows in the value parameter during
the call. ga68 and a68j copy, as the Report requires.
\end{implnote}

Since an array variable has fixed bounds, an assignment must supply an
array with the same bounds. A mismatch is detected when the program runs:

\example{ch04/bounds-mismatch}

\section{Multi-dimensional arrays}\label{sec-matrices}

An array can have several dimensions; the bounds of each dimension are
separated by commas, and so are the indices:

\example{ch04/matrix}

\code|1 UPB m| is the upper bound of the first dimension and
\code|2 UPB m| that of the second; plain \code|UPB m| means
\code|1 UPB m|. A literal for a two-dimensional array is a list of rows,
each of which is a list.

An array of arrays is something different: \code|[] [] int| is an array
whose elements are themselves arrays, and those may have different
lengths, as in \code|((1), (2, 3), (4, 5, 6))|. Its elements are selected
with two pairs of brackets, \code|x[3][2]|. A two-dimensional array, by
contrast, is always rectangular.

\section{Slices}\label{sec-slices}

An index range in place of an index selects a part of an array, a
\term{slice}:

\example{ch04/slices}

The slice \code|a[2:4]| consists of the elements 2 to 4 of \code|a|, but
its own bounds are 1 to 3: a range is renumbered from 1 unless a new lower
bound is given after \code|@|, as in \code|a[2:4 @ 2]|. Either end of the
range can be left out; it then defaults to the corresponding bound of the
array. The Report calls an index range a \rr{trimmer}.

A slice of an array variable is not a copy: it refers to the elements of
the original array, so an assignment to the slice changes the array, as
the last line shows. In the terms of Section~\ref{sec-references}, the
slice of a reference to an array is a reference to part of the array.
An assignment between overlapping parts of an array is safe, since the
right-hand side is a value that is complete before anything is stored:

\example{ch04/shift}

\noindent shifts the elements one place to the right, which in C would
need \texttt{memmove} rather than \texttt{memcpy}.

In a multi-dimensional array, an index can be given for some dimensions
and left out for others. Leaving out the second index of a matrix selects
a row, leaving out the first selects a column:

\example{ch04/matrix-slices}

A dimension whose index is left out entirely keeps its bounds; only a
range with a colon is renumbered:

\example{ch04/slice-bounds}

\noindent For a matrix whose bounds do not start at 1, a loop over a row
\code|m[i, ]| must therefore run from \code|LWB m[i, ]|, not from 1.

A row or an element that a loop body uses several times can be given a
name of its own. Inside a loop over \code|i|, the declaration
\code|ref [] int row = m[i, ]| makes \code|row[j]| the same cell as
\code|m[i, j]|, so assignments to it change the matrix. For an element
that is only read, \code|int x = v[i]| names its current value; this is
a copy, which later assignments to \code|v[i]| do not change.

\section{Strings are arrays}\label{sec-strings}

The type \code|string| is an array of characters. Strings normally have
the lower bound 1; a string variable takes the bounds of whatever is
assigned to it, so a slice with \code|@| can give it another, and loops
should use \code|LWB s|. Everything said about arrays applies to it:
\code|UPB s| is normally its length,
\code|s[i]| its \code|i|-th character, and \code|s[2:4]| a part of it:

\example{ch04/strings}

Strings are compared with \code|=| and \code|<| in the order of their
character codes, which is alphabetical for letters of the same case
(in ASCII, every capital letter comes before every small letter), and
joined with \code|+|. A string variable can hold strings of different
lengths, which an ordinary array cannot; Chapter~\ref{ch-advanced} explains how.

\section{Records}\label{sec-records}

A \term{record} groups values of possibly different types into one value.
Its type, written with \code|struct|, lists the types and names of its
\term{fields}:

\example{ch04/records}

The record type of \code|p| has three fields: \code|name|, a
\code|string|, \code|born|, an \code|int|, and \code|alive|, a
\code|bool|. A field is selected with \code|of|: \code|name of p| is what
most languages write as \texttt{p.name}. If \code|p| is a variable,
\code|born of p| is a reference to the field, so the field can be
assigned to on its own. The Report calls a record a \rr{structure}, a field
name a \rr{field selector}, and \code|name of p| a \rr{selection}.

A record value can be written as a list of field values in parentheses,
in the order of the fields (\rr{structure display}). Such a literal needs at
least two fields; a record with a single field must be built by
assigning to the field.

Selection with \code|of| binds less tightly than indexing:
\code|name of class[i]| means \code|name of (class[i])|, the name of the
\code|i|-th element of an array of records. If a field is itself an array,
parentheses are needed to index it: \code|(marks of class[i])[k]|.
The same holds for a field that is a function: \code|(f of s) (x)|
calls it, while \code|f of s (x)| selects the field \code|f| of
\code|s (x)| and is rejected, with a message that \code|s| must yield a function.

\section{Type declarations}\label{sec-modes}

Writing out a record type at each use would be tedious and error-prone.
A \term{type declaration} (\rr{mode declaration}) gives a type a name:

\example{ch04/modes}

The names of declared types begin with a capital letter
(Section~\ref{sec-stropping}). A declared type can be used wherever a type
is expected, including in other type declarations: a \code|Circle|
contains a \code|Point|. The field \code|x of center of c| is the
\code|x| of the \code|center| of \code|c|.

A type declaration does not create a new, distinct type; it gives a name
to a type that is determined by its structure. Two record types with the
same field types and field names, in the same order, are the same type,
whatever they are called:

\example{ch04/equivalence}

This \term{structural equivalence} differs from the \term{nominal typing}
of C, Java, Rust and most current languages, where two separately
declared \texttt{struct}s are different types even if their fields are
identical. In Algol 68, a \code|Meters| and a \code|Seconds| as declared
above are interchangeable, and the type system does not keep them apart.
Chapter~\ref{ch-advanced} returns to the question of when two types are the same.

A type declaration can also name an array type, such as
\code|mode Vector = [1:3] real|. It must then give the bounds: a type
without bounds, such as \code|[] int|, is allowed for a parameter or a
constant, which take their bounds from a value, but not in a type
declaration, since a type declaration may be used to declare variables.

\section{Arrays of records}\label{sec-multiple-selection}

When the fields of an array of records are selected with \code|of|, the
result is an array of fields:

\example{ch04/multiple-selection}

\code|x of path| is an array of the four \code|x| values. If \code|path|
is a variable, it is even a reference to them, so that
\code|x of path := (1, 2, 3, 4)| assigns to the \code|x| field of every
element. The Report calls this \rr{multiple selection}. Few current
languages have it; the closest analog is selecting a column of a table
in a data-frame library.

\section{A single value as an array}\label{sec-rowing}

Where an array is required and a single value is given, the value is
converted into an array with one element, the bounds 1 to 1:

\example{ch04/rowing}

This implicit conversion is called \term{rowing} (\rr{rowing}) in the
Report, since an array is a \rr{row}. It is why \code|print (newline)|
needs only one pair of parentheses: \code|print| expects a list, and the
single item \code|newline| is turned into a list of one. It is also why
\code|y of ps := 0|, for an array \code|ps| of three points, fails when
the program runs: \code|0| becomes an array with one element, which does
not fit the three \code|y| fields.

\section{Printing arrays and records}\label{sec-print-data}

\code|print| accepts arrays of numbers, truth values and characters,
and prints their elements one after the other, as the examples of this
chapter show. A matrix is printed row by row, without line breaks
between the rows. The Report's \code|print| also accepts records,
which it prints field by field, as long as they contain no references,
functions or unions. A \code|string| field is printed as its characters,
and an array of strings as its strings one after the other.

\begin{implnote}
\code|Genie_Transput| does not accept records in \code|print|; their
fields must be printed one by one. Algol 68 Genie accepts them. All three
implementations print arrays of strings.
\end{implnote}

\section{Two complete programs}\label{sec-data-complete}

The first program sorts an array with \term{insertion sort}: each element
in turn is moved to the left until the elements before it are not
greater. The constant array \code|data| is copied into a variable with
the same bounds, which the program then rearranges.

\example{ch04/sort}

The condition of the inner loop uses \code|andth| (Section~\ref{sec-logic}):
when \code|j| has passed the lower bound, \code|a[j]| must not be
evaluated, since the index would be out of range. With \code|AND|, both
operands would be evaluated and the program would stop with an index
error.

The second program keeps a table of records, each with a name and three
marks, and finds the student with the best average:

\example{ch04/students}

The record type \code|Student| has an array field with the bounds 1 to
3, and the array \code|class| takes its bounds from the literal. Note the
parentheses in \code|(marks of class[i])[k]|, explained in
Section~\ref{sec-records}. The
bounds in a type apply only where a variable is created: they are
computed anew for each variable, and a constant of type \code|Student|
takes the bounds of its value. A literal with four marks would be
accepted in \code|class|, and the error would appear only when such a
record is assigned to a \code|Student| variable.

\begin{implnote}
ga68 and a68j do not check the bounds of an array field when a record is
assigned: a \code|Student| variable then holds four marks. Algol 68 Genie
reports the error.
\end{implnote}

\section*{Exercises}
\addcontentsline{toc}{section}{Exercises}

The \hyperref[ans-ch4]{answers} are at the end of the book.

\begin{enumerate}
\item Declare an array that holds one \code|real| for each year from 1900
  to 2000, and write a loop that sets all elements to 0.
\item Given \code|[] int a = (5, 6, 7, 8, 9)|, what are \code|a[2:3]|,
  \code|LWB a[2:3]|, \code|UPB a[3:]| and \code|a[2:3 @ 0][0]|?
\item Given \code|[1:3, 1:4] int m|, what are the types and bounds of
  \code|m[2, ]|, \code|m[, 3]| and \code|m[2:3, 2]|?
\item After \code|[1:3] int a := (1, 2, 3), b := (4, 5, 6)|, what does
  \code|a := b; b[1] := 0| leave in \code|a|?
\item Declare a type \code|Date| with fields \code|day|, \code|month| and
  \code|year|, and a constant for 3 November 1975.
\item Why does \code|name of class[1][1]| not give the first character of
  the first name in an array \code|class| of records with a
  \code|string| field \code|name|? How must it be written?
\item Given \code|[1:10] Point ps| with \code|mode Point = struct (real x, y)|,
  write a single assignment that sets the \code|y| field of all ten points
  to 0.
\item For \code|[0:9, 2:11] int m|, what are the bounds of \code|m[1, ]|,
  \code|m[1, 3:5]| and \code|m[ , 2]|?
\item What does \code|a[2:5] := a[1:4]| leave in
  \code|[6] int a := (1, 2, 3, 4, 5, 6)|?
\end{enumerate}

\section*{Programming problems}
\addcontentsline{toc}{section}{Programming problems}

\begin{enumerate}
\item Read ten numbers and print them in reverse order.
\item Read a matrix of $3 \times 3$ integers and print its transpose and
  the sum of its diagonal.
\item Multiply two matrices whose sizes are read from the input.
\item Use the sieve of Eratosthenes, an array of \code|bool|, to print
  the prime numbers below 1000.
\item Keep a table of records of the planets (name, distance from the sun,
  diameter) and print them sorted by diameter.
\item \textbf{Binary search.} Write a binary search for a value in a
  sorted array that works for any bounds, using \code|LWB| and
  \code|UPB|. Test it on arrays with the bounds 1 to 10, 0 to 9 and
  $-5$ to 4.
\item \textbf{Merging.} Merge two sorted arrays into one sorted array,
  so that a value that occurs in both appears only once.
\item \textbf{Rotation.} Rotate a square matrix by 90 degrees using row
  and column slices (\code|r[ , n + 1 - i] := m[i, ]|), and rotate an
  array by $r$ places with slice assignments only.
\item \textbf{Game of Life.} Play Conway's Game of Life on a
  \code|[0:n+1, 0:m+1] bool| board whose border stays empty, so that the
  loop over the neighbors needs no special cases. Compute each generation
  into a second board and assign it to the first.
\item \textbf{Generalized Fibonacci.} In the Fibonacci numbers of
  order $m$, each number is the sum of the $m$ before it, starting with
  $m - 1$ zeros and a one. Print the first 30 for a given $m$, keeping
  only the last $m$ numbers in an array \code|[0 : m - 1] int| indexed
  with \code|MOD|.
\item \textbf{League table.} Read the results of the games of a
  league, keep a table of records (team, games, points, goals for and
  against), and print it sorted by points, then by goal difference.
\end{enumerate}
